%\section{Taking turns speaking is the smartest thing to do}
\section{Taking turns speaking (be polite)}
Lets recap. Our problem is that, without a headset, we have that the
signal that we send is
\begin{equation}
    \mathbf{m} = \tilde{\mathbf{n}} + \hat{\mathbf{s}}
\end{equation}
where $\tilde{\mathbf{n}}$ is the version of our voice captured by our
mic, and $\hat{\mathbf{s}}$ is a aproximated-and-multy-echo version of
$\mathbf{s}$, the signal played by our speaker(s). And if we were able
to make $\hat{\mathbf{s}}=\mathbf{0}$, our feedback problem would
vanish, ... at least, theoretically.

But wait ..., as a polite person, if I don't speak when I am listening
to (my interlocutor), I could assume that
$\hat{\mathbf{s}}=\mathbf{0}$ when I am speaking because my
interlocutor should do the same (be silent when I speak)! But, if this
is not true? A solution is: if I am speaking, can attenuate $\mathbf{s}$
(the signal played by my speaker).

Summarizing: you should be able to control the volume of $s$ (the
analog signal that comes out from my speaker) and cut it down (or
even off, if necessary) when you are speaking or your mic is recording
a signal $m$ with enough energy.

\begin{comment}
Therefore, this solution (decrease $s(t)$ when we are speaking, like
the previous one) does need some technique to detect when I'm
speaking. The advantage here is that we only need to switch off (or at
least cut down enough the) volume of our speaker when I speak,
something that does not consume computational resources. Notice that
the ``voice detector'' algorithm can run in a separate thread/process.
\end{comment}
